\chapter{How to use this guide}
\label{ch:start}

\section{Who this guide is for}

You built EcoExplain-VLM over about four months, and now you have to explain it in a report that your professor will read closely. This guide assumes you are starting from zero. It explains every idea the project uses, shows exactly what was done at each step, gives every result with the file it came from, and then helps you check and rewrite your own draft.

The guide follows your own workflow figure (Figure~\ref{fig:workflow-yours} in Chapter~\ref{ch:bigpicture}). That figure has five stages, numbered 01 to 05. Part~II of the guide has one chapter for each of the first four stages, Part~III explains the models and the scoring you need for stage~05, and Part~IV covers the three evaluation tracks of stage~05 one by one.

\section{What is in each part}

\begin{description}[style=nextline,leftmargin=1.2cm]
  \item[Part I. Start here (Chapters 1--3)] How to read the guide; the whole project on a few pages (the workflow figure explained box by box, the order in which the work happened, the numbers to know by heart); and every result in one table.
  \item[Part II. The workflow, stage by stage (Chapters 4--7)] Stage~01, the five data sources, with the remote-sensing ideas behind them. Stage~02, Earth Engine and image processing. Stage~03, how the labels and the evidence cards were made. Stage~04, the frozen dataset and the train/development/test split.
  \item[Part III. Models and measurement (Chapters 8--9)] What a vision--language model is, the four models used and exactly what each one received, LoRA fine-tuning, and every metric used to score classifications and explanations.
  \item[Part IV. The three evaluation tracks (Chapters 10--12)] Track~A, zero-shot classification. Track~B, class LoRA. Track~C, explanations (free-form, structured, and the explanation LoRA). Each chapter gives the results, how to read them, and what they do not show.
  \item[Part V. Putting it together (Chapter 13)] The answer to the research question, the hypotheses, the objectives, the limitations and the future work.
  \item[Part VI. From understanding to your report (Chapters 14--17)] A check of every figure you have; a section-by-section review of your current draft; a writing plan with word budgets and the references that fit each section; and the questions your professor is likely to ask.
  \item[Appendices] Every key number with its source file, a map of the project folders, a glossary, and the answers to the self-check questions.
\end{description}

\section{A reading plan}

You do not have to read the guide in one go. A plan that works in about a week:

\begin{table}[H]
\centering
\small
\begin{tabular}{lL{6.4cm}L{6.0cm}}
\toprule
Day & Read & Then do this \\
\midrule
1 & Chapters 1--3 & Draw the workflow figure from memory, with the key numbers. \\
2 & Chapters 4--5 (stages 01--02) & Explain NDVI, SAR backscatter and median compositing aloud in two minutes each. \\
3 & Chapters 6--7 (stages 03--04) & Open \file{dataset.jsonl}, find one test patch, and read its evidence card. \\
4 & Chapters 8--9 & Write the LoRA equation and the Wilson interval from memory. \\
5 & Chapters 10--12 & Re-create the main results table on paper and say what each row means. \\
6 & Chapters 13--14 & List the limitations in your own words; decide which figures go in the report. \\
7 & Chapters 15--17 & Go through the draft review and fix your draft, section by section. \\
\bottomrule
\end{tabular}
\end{table}

Each chapter ends with a ``Check yourself'' box. Answer the questions without looking back, then compare with Appendix~\ref{app:answers}. If you can answer them, you can defend that part of the report.

\section{The boxes}

Seven kinds of box appear in the guide.

\begin{plainwords}
A short explanation without jargon.
\end{plainwords}
\begin{whybox}
The reason a choice was made, and what the alternatives were.
\end{whybox}
\begin{watchout}
A problem, inconsistency or trap. These matter most for an honest report.
\end{watchout}
\begin{foryourreport}
What this means for the report you will write. These boxes list what to include. They never contain finished report sentences: the sentences must be yours.
\end{foryourreport}
\begin{wherecode}
The files, functions and notebook cells that implement the idea, so you can check it.
\end{wherecode}
\begin{newinfo}
Something that was added or corrected since an earlier version of the guide (version~1: 6 October; version~2: 19 October 2026). Read these boxes first if you know an earlier version.
\end{newinfo}
\begin{selfcheck}
Questions to test your understanding at the end of each chapter.
\end{selfcheck}

\section{Conventions}

\begin{itemize}
  \item File paths are relative to the project folder \file{ecoexplain_vlm_dataset/} unless they start with \file{ada_review/} (files copied back from the IIIT Ada cluster), \file{report/} or \file{study_guide_v2/} (folders next to the project folder).
  \item Percentages are on the 55 test patches unless stated otherwise.
  \item ``Gold'' or ``label'' always means \texttt{candidate\_class}, the automatic label used for all scoring.
  \item \textbf{Phases} are the order in which the work happened (A, B, C, D). \textbf{Tracks} are the three experiments in stage~05 of your workflow figure (A, B, C). The letters do not mean the same thing; Table~\ref{tab:phasetrack} in Chapter~\ref{ch:bigpicture} maps one onto the other.
\end{itemize}

\section{What changed in version 3}

Version~2 (19 October) added the explanation LoRA, the exact input of each model, a check of your figures and a review of your draft. On 19--20 October you then fixed several protocol problems in the code and re-ran the key experiments. Version~3 brings those results in and corrects what they overturn.

\begin{newinfo}
\begin{enumerate}
  \item \textbf{One chapter with every result} (Chapter~\ref{ch:results}), recomputed from the prediction files by \file{build_results_master.py}. Its table is generated, so it cannot drift from the files.
  \item \textbf{Leak-free zero-shot results} (Chapter~\ref{ch:trackA}). Without the summary line Qwen2-VL-2B scores 32.7\,\% (\texttt{evidence\_only}) and 32.7\,\% (\texttt{compact}, answering ``stable'' for every patch); LLaVA-Interleave-0.5B scores 43.6\,\% (\texttt{compact}). PaliGemma could not be loaded and BLIP-2 gave no parsable answer.
  \item \textbf{A correction of version~2.} Version~2 said Qwen ``did not copy the label'' because it followed the card's trend words. Those trend words appear \emph{only} in the summary line, next to \texttt{candidate\_class}. Qwen did not copy the label token, but it did depend on that line; without it, it falls to chance (Chapter~\ref{ch:trackA}).
  \item \textbf{The balanced, leak-free class LoRA} (Ada job 29575; Chapter~\ref{ch:trackB}): 16 patches (6 decline, 5 stable, 5 increase), label-free card at training and test, 36.4\,\%. It answers ``stable'' for 53 of 55 patches. The old 67.3\,\% run is superseded.
  \item \textbf{Code changes} you made (Chapters~\ref{ch:stage04} and~\ref{ch:models}): the \texttt{compact} text mode that gives every model the same table, the parser that ignores echoed prompts, \texttt{--balanced-train}, and a spatial block split (266/46/54).
  \item \textbf{GitHub.} The code, the new results and a README section (4.7) are on the \texttt{main} branch of \texttt{Rama-Vaibhav/EcoExplain-VLM} (Appendix~\ref{app:files}).
  \item \textbf{The draft review} (Chapter~\ref{ch:draft}) now says which issues you report as fixed and adds the checks the new results require.
\end{enumerate}
\end{newinfo}

The version~2 changes (Phase~D, unequal model inputs, the first-16-by-ID training sets, the real LLaVA checkpoint, the raster-based card text, the two score files) are all still in the guide where they belong.

\section{A promise about numbers}

Every number in this guide was read from a project file or recomputed from one, and the source is named. Where the project does not record something (for example, why the candidate pool has an empty strip), the guide says so instead of guessing. When you write your report, keep to the same rule: if you cannot point to the file a number came from, do not use the number.
